\chapter{States, bases and exact diagonalization}
\setwlchapter{bases}
\section{Fix the mode order first}
Let $|n_1\ldots n_M\rangle=(a^\dagger_1)^{n_1}\cdots
(a^\dagger_M)^{n_M}|0\rangle$. An annihilator acting on slot $j$
picks up $(-1)^{\sum_{i<j}n_i}$ before it reaches that slot. Neither
the binary occupation vector nor the site labels alone determine this
sign: the associated \fn{BASIS} is part of the representation. Keep
the mode order when moving a matrix into another numerical code.
An operator-form state instead retains the creators explicitly and
can be inspected without decoding bit positions.
\begin{wl}
snegfermionoperators[d];
makebasis[{d[]}];
{BASIS, ap[d[AN, DO], vc[1, 1]]}
\end{wl}

\section{Block labels versus degeneracies}
Use charge and spin projection if $[H,N]=[H,S^z]=0$. Use total spin
only if $H$ commutes with all components. A transverse field can break
$S^z$; a pairing term breaks $N$ but may preserve spin and fermion
parity. Projecting onto a non-invariant block loses couplings rather
than proving that they vanish. One diagnostic for a projector $P$ is
$(1-P)HP=0$. A second is to compare the reconstructed full spectrum
for a cluster small enough that symmetries are unnecessary.
\casestudy{multiplets}{Counting states in a highest-weight basis}
\casestudy{reflection}{Separating mirror-even and mirror-odd states}
Reflection exchanges sites, not just occupation numbers. A permutation
of creators must be reordered into the reference order. This can supply
a minus sign even when the \emph{pattern} of occupied sites is invariant.
Likewise particle--hole conjugation maps $Q$ to $-Q$ and interchanges
local double occupation and the vacuum. On a bipartite lattice at
particle--hole symmetry one can use a continuous charge isospin
instead of just pairing states of opposite charge.
\casestudy{isospin}{Spin and charge multiplets on a bipartite dimer}

\section{Matrix elements, eigenstates and thermal traces}
In each invariant sector the matrix element is $\langle i|H|j\rangle$,
with the ket indexed by the column. Wolfram Language
\fn{Eigensystem} returns eigenvectors as a list of vectors. For a
numerical normalized eigenvector $v$ compute an expectation as
$v^\dagger A v$, not $v^T A v$. For an operator changing quantum
numbers, such as $c^\dagger$, a matrix with both indices in the same
sector may be identically zero. Use \fn{makeallmatricesbzvc} to obtain
the rectangular blocks between sectors.
\casestudy{thermal}{A partition function from reduced spin blocks}
\begin{exercise}[Weight of a multiplet]\label{ex:multiplet}
The two-spin Heisenberg dimer has one singlet at $-3J/4$ and one
triplet at $J/4$. Derive $\langle\vS_1\cdot\vS_2\rangle_T$ and
check both temperature limits for $J>0$.
\end{exercise}
